\documentclass[14pt]{extarticle}
% =============================================================================
% Shared Preamble for Worksheets, Homeworks, Exams & Lesson Notes
% (scripts/tex_preamble.tex)
% =============================================================================

\newif\ifdetailed
\ifdefined\setdetailed
  \setdetailed
\fi

\newif\ifwithsols
\ifdefined\setwithsols
  \setwithsols
\fi

\newif\ifwhiteblack
\ifdefined\setwhiteblack
  \setwhiteblack
\fi

% Document-class specific packages (loaded before common_preamble.tex / polyglossia)
\usepackage{fancyhdr}
\usepackage{enumitem}
\usepackage{titlesec}
\usepackage[bottom]{footmisc}
\usepackage{titling}
\usepackage{longtable} % page-breakable tables
\usepackage{float}
\usepackage[onehalfspacing]{setspace}

\ifwhiteblack
  \usepackage[a4paper, landscape, top=1.25cm, bottom=1.25cm, left=1.25cm, right=1.25cm, headheight=1.5cm, headsep=0.5cm]{geometry}
\else
  \usepackage[a4paper, top=2.4cm, bottom=2.4cm, left=2.2cm, right=2.2cm, headheight=1.4cm, headsep=0.4cm, footskip=1.0cm]{geometry}
\fi

% Shared common packages, fonts, polyglossia, callout boxes, code helpers, date & header macros
\IfFileExists{scripts/common_preamble.tex}{%
  % =============================================================================
% Shared Common Preamble (scripts/common_preamble.tex)
% Included by:
%   - scripts/tex_preamble.tex    (worksheets, homeworks, exams, lesson notes)
%   - scripts/book_preamble.tex   (two-sided Arabic book)
%   - scripts/beamer_preamble.tex (Beamer presentation slides)
% =============================================================================

\usepackage{fontspec}
\usepackage{amsmath}
\usepackage{amssymb}
\usepackage{xcolor}
\usepackage{graphicx}
\usepackage{listings}
\usepackage[most]{tcolorbox}
\usepackage{pifont}
\usepackage{minted}
\usepackage{etoolbox}
\usepackage{array}
\usepackage{booktabs}
\usepackage{tikz}
\usepackage{datetime2}

\usetikzlibrary{automata, positioning, arrows.meta, arrows, shapes.geometric, calc}

% Beamer already loads hyperref internally; load it here for article/book classes
% right before polyglossia (bidi).
\makeatletter
\@ifclassloaded{beamer}{}{%
  \usepackage{hyperref}%
  \hypersetup{hidelinks}%
  \providecommand{\LTcaptype}{table}%
}
\makeatother

% Load polyglossia (and bidi) AFTER graphics, tables, boxes, code, and hyperref
\usepackage{polyglossia}
\setmainlanguage[numerals=western]{arabic}
\setotherlanguage{english}
\newfontfamily\arabicfont[Script=Arabic]{Calibri}
\newfontfamily\arabicfontsf[Script=Arabic]{Calibri}
\newfontfamily\arabicfonttt[Script=Arabic]{Courier New}
\newfontfamily\emoji[BoldFont=*,ItalicFont=*]{Segoe UI Emoji}

% Allow slight inter-word stretch to avoid Overfull \hbox with inline \textenglish
\emergencystretch=1.5em

% -----------------------------------------------------------------------------
% Default Minted & Callout Box Definitions
% -----------------------------------------------------------------------------
\AtBeginEnvironment{minted}{\setstretch{0.95}}
\setminted{fontsize=\small, breaklines=true, vspace=2pt}

\tcbset{
  arc=2.5pt, outer arc=2.5pt,
  boxsep=1pt, left=6pt, right=6pt, top=3pt, bottom=3pt,
  toptitle=1.5pt, bottomtitle=1.5pt,
  before skip=6pt, after skip=6pt,
  fonttitle=\bfseries,
  fontupper=\setstretch{1.1}\setlength{\parskip}{2pt},
  colback=white, colframe=black, boxrule=0.6pt,
  boxcolor/.style 2 args={colback=#1, colframe=#2}
}

\newtcolorbox{boxDef}[1][]{boxcolor={blue!4!white}{blue!75!black},
  title={{\emoji📘} تعريف\ifx\\#1\\\else ~#1\fi :}}
\newtcolorbox{boxExercise}[1][]{boxcolor={cyan!4!white}{cyan!70!black},
  title={{\emoji🧩} تمرين\ifx\\#1\\\else ~#1\fi :}}
\newtcolorbox{boxExample}[1][]{boxcolor={yellow!6!white}{orange!85!black},
  title={{\emoji📝} مثال\ifx\\#1\\\else ~#1\fi :}}
\newtcolorbox{boxNote}[1][]{boxcolor={gray!7!white}{black!75},
  title={{\emoji✨} ملاحظة\ifx\\#1\\\else ~#1\fi :}}
\newtcolorbox{boxAttention}[1][]{boxcolor={magenta!5!white}{magenta!75!black},
  title={{\emoji🔔} تنبيه\ifx\\#1\\\else ~#1\fi :}}
\newtcolorbox{boxWarning}[1][]{boxcolor={red!4!white}{red!75!black},
  title={{\emoji⚡} ملاحظة هامة\ifx\\#1\\\else ~#1\fi :}}
\newtcolorbox{boxSolution}[1][]{boxcolor={green!4!white}{green!60!black},
  title={{\emoji✅} حل\ifx\\#1\\\else ~#1\fi :}}
\newtcolorbox{boxSymbol}[1][]{boxcolor={purple!4!white}{purple!70!black},
  title={{\emoji🔣} رمز\ifx\\#1\\\else ~#1\fi :}}
\newtcolorbox{boxHint}[1][]{boxcolor={teal!4!white}{teal!65!black},
  title={{\emoji💡} تلميح\ifx\\#1\\\else ~#1\fi :}}

% Compact code boxes: tighter line spacing and small vertical gaps
\tcbset{
  simplecode/.style={
    arc=2pt, outer arc=2pt,
    left=4pt, right=4pt, top=4pt, bottom=4pt,
    fontupper=\setstretch{0.95},
    before skip=6pt, after skip=6pt,
    colback=gray!5, colframe=black!50, boxrule=0.4pt
  }
}
\newenvironment{boxCode}{\begin{tcolorbox}[simplecode]}{\end{tcolorbox}}

% Code cell used inside tables (keeps code left-to-right inside RTL tables)
% Use inside a fixed-width p{...} column: \begin{pycell}\begin{minted}...\end{minted}\end{pycell}
\makeatletter
\newtcolorbox{pycell}{frame empty,colback=white,boxsep=0pt,left=2pt,right=2pt,top=2pt,bottom=2pt,width=\linewidth,before=\noindent,after=,baseline={\tcb@height-\ht\strutbox},fontupper=\setstretch{0.95}\setlength{\topsep}{0pt}\setlength{\partopsep}{0pt}\setlength{\parskip}{0pt}}
\makeatother

\newcolumntype{C}[1]{>{\centering\arraybackslash}p{#1}}

% Dotted date style (DD.MM.YYYY)
\makeatletter
\DTMnewdatestyle{dotted}{%
  \renewcommand{\DTMdisplaydate}[4]{\two@digits{##3}.\two@digits{##2}.##1}%
  \renewcommand{\DTMDisplaydate}{\DTMdisplaydate}%
}
\makeatother
\DTMsetdatestyle{dotted}

% -----------------------------------------------------------------------------
% Customizable Header / Footer Metadata
% -----------------------------------------------------------------------------
\newcommand{\headerleft}{مدرسة التسامح الشاملة}
\newcommand{\headerright}{الأستاذ محمود اغبارية}
\newcommand{\footercenter}{\thepage}

\newcommand{\setheaderleft}[1]{\renewcommand{\headerleft}{#1}}
\newcommand{\setheaderright}[1]{\renewcommand{\headerright}{#1}}
\newcommand{\setfootercenter}[1]{\renewcommand{\footercenter}{#1}}
%
}{%
  \IfFileExists{../../../scripts/common_preamble.tex}{%
    % =============================================================================
% Shared Common Preamble (scripts/common_preamble.tex)
% Included by:
%   - scripts/tex_preamble.tex    (worksheets, homeworks, exams, lesson notes)
%   - scripts/book_preamble.tex   (two-sided Arabic book)
%   - scripts/beamer_preamble.tex (Beamer presentation slides)
% =============================================================================

\usepackage{fontspec}
\usepackage{amsmath}
\usepackage{amssymb}
\usepackage{xcolor}
\usepackage{graphicx}
\usepackage{listings}
\usepackage[most]{tcolorbox}
\usepackage{pifont}
\usepackage{minted}
\usepackage{etoolbox}
\usepackage{array}
\usepackage{booktabs}
\usepackage{tikz}
\usepackage{datetime2}

\usetikzlibrary{automata, positioning, arrows.meta, arrows, shapes.geometric, calc}

% Beamer already loads hyperref internally; load it here for article/book classes
% right before polyglossia (bidi).
\makeatletter
\@ifclassloaded{beamer}{}{%
  \usepackage{hyperref}%
  \hypersetup{hidelinks}%
  \providecommand{\LTcaptype}{table}%
}
\makeatother

% Load polyglossia (and bidi) AFTER graphics, tables, boxes, code, and hyperref
\usepackage{polyglossia}
\setmainlanguage[numerals=western]{arabic}
\setotherlanguage{english}
\newfontfamily\arabicfont[Script=Arabic]{Calibri}
\newfontfamily\arabicfontsf[Script=Arabic]{Calibri}
\newfontfamily\arabicfonttt[Script=Arabic]{Courier New}
\newfontfamily\emoji[BoldFont=*,ItalicFont=*]{Segoe UI Emoji}

% Allow slight inter-word stretch to avoid Overfull \hbox with inline \textenglish
\emergencystretch=1.5em

% -----------------------------------------------------------------------------
% Default Minted & Callout Box Definitions
% -----------------------------------------------------------------------------
\AtBeginEnvironment{minted}{\setstretch{0.95}}
\setminted{fontsize=\small, breaklines=true, vspace=2pt}

\tcbset{
  arc=2.5pt, outer arc=2.5pt,
  boxsep=1pt, left=6pt, right=6pt, top=3pt, bottom=3pt,
  toptitle=1.5pt, bottomtitle=1.5pt,
  before skip=6pt, after skip=6pt,
  fonttitle=\bfseries,
  fontupper=\setstretch{1.1}\setlength{\parskip}{2pt},
  colback=white, colframe=black, boxrule=0.6pt,
  boxcolor/.style 2 args={colback=#1, colframe=#2}
}

\newtcolorbox{boxDef}[1][]{boxcolor={blue!4!white}{blue!75!black},
  title={{\emoji📘} تعريف\ifx\\#1\\\else ~#1\fi :}}
\newtcolorbox{boxExercise}[1][]{boxcolor={cyan!4!white}{cyan!70!black},
  title={{\emoji🧩} تمرين\ifx\\#1\\\else ~#1\fi :}}
\newtcolorbox{boxExample}[1][]{boxcolor={yellow!6!white}{orange!85!black},
  title={{\emoji📝} مثال\ifx\\#1\\\else ~#1\fi :}}
\newtcolorbox{boxNote}[1][]{boxcolor={gray!7!white}{black!75},
  title={{\emoji✨} ملاحظة\ifx\\#1\\\else ~#1\fi :}}
\newtcolorbox{boxAttention}[1][]{boxcolor={magenta!5!white}{magenta!75!black},
  title={{\emoji🔔} تنبيه\ifx\\#1\\\else ~#1\fi :}}
\newtcolorbox{boxWarning}[1][]{boxcolor={red!4!white}{red!75!black},
  title={{\emoji⚡} ملاحظة هامة\ifx\\#1\\\else ~#1\fi :}}
\newtcolorbox{boxSolution}[1][]{boxcolor={green!4!white}{green!60!black},
  title={{\emoji✅} حل\ifx\\#1\\\else ~#1\fi :}}
\newtcolorbox{boxSymbol}[1][]{boxcolor={purple!4!white}{purple!70!black},
  title={{\emoji🔣} رمز\ifx\\#1\\\else ~#1\fi :}}
\newtcolorbox{boxHint}[1][]{boxcolor={teal!4!white}{teal!65!black},
  title={{\emoji💡} تلميح\ifx\\#1\\\else ~#1\fi :}}

% Compact code boxes: tighter line spacing and small vertical gaps
\tcbset{
  simplecode/.style={
    arc=2pt, outer arc=2pt,
    left=4pt, right=4pt, top=4pt, bottom=4pt,
    fontupper=\setstretch{0.95},
    before skip=6pt, after skip=6pt,
    colback=gray!5, colframe=black!50, boxrule=0.4pt
  }
}
\newenvironment{boxCode}{\begin{tcolorbox}[simplecode]}{\end{tcolorbox}}

% Code cell used inside tables (keeps code left-to-right inside RTL tables)
% Use inside a fixed-width p{...} column: \begin{pycell}\begin{minted}...\end{minted}\end{pycell}
\makeatletter
\newtcolorbox{pycell}{frame empty,colback=white,boxsep=0pt,left=2pt,right=2pt,top=2pt,bottom=2pt,width=\linewidth,before=\noindent,after=,baseline={\tcb@height-\ht\strutbox},fontupper=\setstretch{0.95}\setlength{\topsep}{0pt}\setlength{\partopsep}{0pt}\setlength{\parskip}{0pt}}
\makeatother

\newcolumntype{C}[1]{>{\centering\arraybackslash}p{#1}}

% Dotted date style (DD.MM.YYYY)
\makeatletter
\DTMnewdatestyle{dotted}{%
  \renewcommand{\DTMdisplaydate}[4]{\two@digits{##3}.\two@digits{##2}.##1}%
  \renewcommand{\DTMDisplaydate}{\DTMdisplaydate}%
}
\makeatother
\DTMsetdatestyle{dotted}

% -----------------------------------------------------------------------------
% Customizable Header / Footer Metadata
% -----------------------------------------------------------------------------
\newcommand{\headerleft}{مدرسة التسامح الشاملة}
\newcommand{\headerright}{الأستاذ محمود اغبارية}
\newcommand{\footercenter}{\thepage}

\newcommand{\setheaderleft}[1]{\renewcommand{\headerleft}{#1}}
\newcommand{\setheaderright}[1]{\renewcommand{\headerright}{#1}}
\newcommand{\setfootercenter}[1]{\renewcommand{\footercenter}{#1}}
%
  }{%
    \input{../../scripts/common_preamble.tex}%
  }%
}

\pagestyle{fancy}

% Prevent widow/orphan lines
\widowpenalty=10000
\clubpenalty=10000

% Compact section headings and default list spacing so pages do not waste vertical space
\titlespacing*{\section}{0pt}{1.4ex plus 0.4ex minus 0.2ex}{0.7ex plus 0.2ex}
\titlespacing*{\subsection}{0pt}{1.1ex plus 0.3ex minus 0.2ex}{0.5ex plus 0.1ex}
\titlespacing*{\subsubsection}{0pt}{0.9ex plus 0.2ex minus 0.1ex}{0.4ex plus 0.1ex}
\setlist{topsep=4pt, partopsep=0pt, parsep=2pt, itemsep=4pt}

\gappto\captionsarabic{\renewcommand{\contentsname}{جدول المحتويات}}

% Legacy listings settings
\lstset{
  language=[Sharp]C,
  numbers=left,
  stepnumber=1,
  numbersep=8pt,
  frame=single,
  basicstyle=\ttfamily\small
}
\ifwhiteblack
  \lstset{
    keywordstyle=\bfseries,
    stringstyle=\ttfamily,
    commentstyle=\itshape
  }
\else
  \lstset{
    keywordstyle=\color{blue},
    stringstyle=\color{red},
    commentstyle=\color{green!50!black}
  }
\fi

% White-and-black landscape presentation mode overrides
\ifwhiteblack
  \AtBeginDocument{
    \LARGE
    \usemintedstyle{bw}
  }
  \setminted{vspace=2pt}
  \tcbset{
    colback=white, colframe=black, boxrule=1pt,
    boxcolor/.style 2 args={colback=white, colframe=black},
    simplecode/.append style={colback=white, colframe=black, boxrule=1pt}
  }
\fi

% Default to empty author to avoid "LaTeX Warning: No \author given" when omitted
\author{}

% Compact title block
\makeatletter
\renewcommand{\@maketitle}{%
  \begin{center}%
    {\ifwhiteblack\Huge\else\LARGE\fi \bfseries \@title \par}%
    \ifdefempty{\@author}{}{%
      \vskip 0.15em%
      {\large \@author \par}%
    }%
    {\small \hfill \textenglish{\DTMtoday} \par}%
  \end{center}%
  \vspace{-0.3cm}%
}
\makeatother

\apptocmd{\maketitle}{\thispagestyle{fancy}}{}{}

% Apply page headers and footer (using \headerleft, \headerright, \footercenter from common_preamble.tex)
\newcommand{\applyfancyheaders}{%
  \fancyhf{}%
  \fancyhead[L]{\textarabic{\headerleft}}%
  \fancyhead[R]{\textarabic{\headerright}}%
  \fancyfoot[C]{\textarabic{\footercenter}}%
}

\fancypagestyle{plain}{\applyfancyheaders}
\applyfancyheaders

\setcounter{tocdepth}{3} % only section, subsection, and subsubsection in TOC

\newcommand{\importpdfpage}[2]{%
    \newgeometry{left=0cm,right=0cm,top=0cm,bottom=0cm}
    \thispagestyle{empty}
    \includegraphics[page=#2,width=0.95\paperwidth,height=0.95\paperheight,keepaspectratio]{#1}
    \restoregeometry
}

\newcommand{\insertFullImg}[1]{%
    \noindent
    \makebox[\textwidth][c]{
        \includegraphics[width=0.9\paperwidth,keepaspectratio]{#1}%
    }%
}%


\ifwithsols
\title{حل ورقة تمرن $7$ للصف العاشر $10$ - حلقات التكرار \textenglish{while} و \textenglish{for}}
\else
\title{ورقة تمرن $7$ للصف العاشر $10$ - حلقات التكرار \textenglish{while} و \textenglish{for}}
\fi

\begin{document}

\maketitle
\thispagestyle{fancy}

% ==============================================================
% القسم الأول: تمهيد وتتبّع وإكمال كود في حلقة while
% ==============================================================
\section*{القسم الأوّل: تمهيد وتتبّع وإكمال كود في حلقة \textenglish{while}}

\begin{enumerate}[itemsep=2em]

\item تَتبَّع مقطع البرنامج التالي وأجب عن الأسئلة التي تليه:
\begin{boxCode}
\begin{english}
\begin{minted}{csharp}
int x = 2;
while (x <= 10)
{
    Console.WriteLine(x);
    x = x + 3;
}
Console.WriteLine("End: " + x);
\end{minted}
\end{english}
\end{boxCode}
\begin{enumerate}
    \item ما هي الأعداد التي يطبعها البرنامج داخل الحلقة؟
    \item ما هي قيمة $x$ النهائية التي جعلت شرط الحلقة يصبح \textenglish{false} وتتوقف الحلقة؟
\end{enumerate}
\ifwithsols
\begin{boxSolution}
\begin{enumerate}
    \item يطبع داخل الحلقة الأعداد: $2$، ثم $5$، ثم $8$.
    \item قيمة $x$ عند التوقف هي $11$ (لأنّ $11 \le 10$ غير متحقّق)، فيطبع في النهاية: \textenglish{End: 11}.
\end{enumerate}
\end{boxSolution}
\fi
\filbreak

\item تَتبَّع مقطع البرنامج التالي بواسطة جدول متابعة واكتب مخرجاته:
\begin{boxCode}
\begin{english}
\begin{minted}{csharp}
int num = 70;
while (num >= 61)
{
    if (num % 2 == 0)
    {
        Console.WriteLine(num);
    }
    num--;
}
\end{minted}
\end{english}
\end{boxCode}
\ifwithsols
\begin{boxSolution}
يطبع البرنامج الأعداد الزوجية من $70$ نزولاً إلى $62$:
\begin{english}
\begin{minted}{text}
70
68
66
64
62
\end{minted}
\end{english}
\end{boxSolution}
\fi
\filbreak

\item تَتبَّع مقطع البرنامج التالي:
\begin{boxCode}
\begin{english}
\begin{minted}{csharp}
int n = int.Parse(Console.ReadLine());
while (n != 0)
{
    Console.WriteLine(n * 2);
    n = int.Parse(Console.ReadLine());
}
Console.WriteLine("Finished");
\end{minted}
\end{english}
\end{boxCode}
\begin{enumerate}
    \item إذا أدخل المستخدم الأعداد: $4$ ثم $7$ ثم $-3$ ثم $0$، ماذا يطبع البرنامج؟
    \item إذا كان أوّل عدد يُدخله المستخدم هو $0$، ماذا يطبع البرنامج؟
    \item لماذا نحتاج السطر \textenglish{n = int.Parse(Console.ReadLine());} في نهاية جسم الحلقة؟ ماذا سيحدث لو حذفناه؟
\end{enumerate}
\ifwithsols
\begin{boxSolution}
\begin{enumerate}
    \item يطبع: $8$، ثم $14$، ثم $-6$، ثم \textenglish{Finished} (لاحظ أنّ الـ $0$ لا يُضرب بـ $2$ لأنّ الحلقة تتوقف عنده).
    \item يطبع \textenglish{Finished} فقط، لأنّ الشرط \textenglish{n != 0} يكون \textenglish{false} منذ البداية فلا يدخل إلى جسم الحلقة أبدًا.
    \item نحتاجه لقراءة العدد التالي في كل دورة؛ لو حذفناه ستبقى قيمة $n$ الأولى ثابتة إلى الأبد ويدخل البرنامج في حلقة لا نهائية.
\end{enumerate}
\end{boxSolution}
\fi
\filbreak

\item تَتبَّع مقطع البرنامج التالي:
\begin{boxCode}
\begin{english}
\begin{minted}{csharp}
int sum = 0;
int num = int.Parse(Console.ReadLine());
while (num != -1)
{
    sum = sum + num;
    num = int.Parse(Console.ReadLine());
}
Console.WriteLine("Sum = " + sum);
\end{minted}
\end{english}
\end{boxCode}
إذا أدخل المستخدم الأعداد: $10$ ثم $5$ ثم $20$ ثم $-1$:
\begin{enumerate}
    \item أكمل جدول المتابعة التالي:
    \begin{center}
    \begin{english}
    \renewcommand{\arraystretch}{1.4}
    \begin{tabular}{|C{2.5cm}|C{3.5cm}|C{2.5cm}|C{3.5cm}|}
    \hline
    \textbf{num} & \textbf{num != -1} & \textbf{sum} & \textarabic{\textbf{المخرجات}} \\
    \hline
     & & $0$ & \\ \hline
    $10$ & \textenglish{true} & & \\ \hline
    $5$ & & & \\ \hline
    $20$ & & & \\ \hline
    $-1$ & & & \\ \hline
    \end{tabular}
    \end{english}
    \end{center}
    \item هل جُمِع العدد $-1$ مع المتغير \textenglish{sum}؟ لماذا؟
\end{enumerate}
\ifwithsols
\begin{boxSolution}
\begin{enumerate}
    \item جدول المتابعة:
    \begin{center}
    \begin{english}
    \renewcommand{\arraystretch}{1.4}
    \begin{tabular}{|C{2.5cm}|C{3.5cm}|C{2.5cm}|C{3.5cm}|}
    \hline
    \textbf{num} & \textbf{num != -1} & \textbf{sum} & \textarabic{\textbf{المخرجات}} \\
    \hline
     & & $0$ & \\ \hline
    $10$ & \textenglish{true} & $10$ & \\ \hline
    $5$ & \textenglish{true} & $15$ & \\ \hline
    $20$ & \textenglish{true} & $35$ & \\ \hline
    $-1$ & \textenglish{false} & $35$ & \textenglish{Sum = 35} \\ \hline
    \end{tabular}
    \end{english}
    \end{center}
    \item كلا، لم يُجمع العدد $-1$ لأنّنا قرأناه في نهاية الدورة الثالثة، وعندما فُحص الشرط \textenglish{num != -1} أعاد \textenglish{false} فخرجت الحلقة فورًا قبل تنفيذ \textenglish{sum = sum + num}.
\end{enumerate}
\end{boxSolution}
\fi
\filbreak

\item \textbf{تدريب على كتابة شرط \textenglish{while}:} \\
تذكّر أنّنا نكتب داخل \textenglish{while (...)} \textbf{شرط الاستمرار} (متى تستمرّ الحلقة؟)، وهو \textbf{عكس شرط التوقف}. \\
لكل حالة في الجدول التالي، اكتب الشرط المناسب داخل \textenglish{while (...)}:

{\centering
\renewcommand{\arraystretch}{1.6}
\begin{tabular}{|>{\raggedleft\arraybackslash}p{0.62\textwidth}|>{\centering\arraybackslash}p{0.32\textwidth}|}
\hline
\textbf{المطلوب في السؤال (متى تتوقف الحلقة؟)} & \textbf{شرط الاستمرار في \textenglish{while}} \\
\hline
$1$. تتوقف الحلقة عندما يُدخل المستخدم العدد $0$ في المتغير $n$. & \textenglish{while (n != 0)} \\
\hline
$2$. تتوقف الحلقة عندما يُدخل المستخدم العدد $-1$ في المتغير $x$. & \textenglish{while (............)} \\
\hline
$3$. تتوقف الحلقة عندما يُدخل المستخدم عددًا \textbf{سالبًا} في المتغير \textenglish{num} (أي تستمر طالما العدد غير سالب). & \textenglish{while (............)} \\
\hline
$4$. تتوقف الحلقة عندما يُدخل المستخدم كلمة المرور الصحيحة $1234$ في المتغير \textenglish{pass}. & \textenglish{while (............)} \\
\hline
$5$. تتوقف الحلقة عندما يصبح المجموع \textenglish{sum} أكبر تمامًا من $100$. & \textenglish{while (............)} \\
\hline
$6$. تتوقف الحلقة عندما يُدخل المستخدم في $n$ عددًا يقبل القسمة على $7$ بدون باقٍ. & \textenglish{while (............)} \\
\hline
\end{tabular}
\par}
\ifwithsols
\begin{boxSolution}
\begin{enumerate}[nosep]
    \item \textenglish{while (n != 0)}
    \item \textenglish{while (x != -1)}
    \item \textenglish{while (num >= 0)}
    \item \textenglish{while (pass != 1234)}
    \item \textenglish{while (sum <= 100)}
    \item \textenglish{while (n \% 7 != 0)}
\end{enumerate}
\end{boxSolution}
\fi
\filbreak

\item \textbf{المستوى $1$ — أكمل شرط الحلقة فقط:}
\begin{enumerate}
    \item البرنامج التالي يستقبل أعدادًا صحيحة ويطبع كل عدد يدخله المستخدم، ويتوقف عندما يُدخل المستخدم العدد $0$. أكمل الشرط الناقص:
\begin{boxCode}
\begin{english}
\begin{minted}{csharp}
int n = int.Parse(Console.ReadLine());
while ( ____________ )
{
    Console.WriteLine(n);
    n = int.Parse(Console.ReadLine());
}
\end{minted}
\end{english}
\end{boxCode}

    \item البرنامج التالي يطلب من المستخدم إدخال عدد موجب (أكبر من $0$). طالما أدخل المستخدم عددًا أصغر أو يساوي $0$، يطبع البرنامج رسالة خطأ ويطلب عددًا جديدًا:
\begin{boxCode}
\begin{english}
\begin{minted}{csharp}
int x = int.Parse(Console.ReadLine());
while ( ____________ )
{
    Console.WriteLine("Error! Enter a positive number:");
    x = int.Parse(Console.ReadLine());
}
Console.WriteLine("Thank you! You entered: " + x);
\end{minted}
\end{english}
\end{boxCode}
\end{enumerate}
\ifwithsols
\begin{boxSolution}
\begin{enumerate}[nosep]
    \item الشرط هو: \textenglish{n != 0}
    \item الشرط هو: \textenglish{x <= 0}
\end{enumerate}
\end{boxSolution}
\fi
\filbreak

\item \textbf{المستوى $2$ — أكمل الشرط وسطْر التحديث داخل الحلقة:}
\begin{enumerate}
    \item البرنامج التالي يطلب من المستخدم إدخال عدد يقبل القسمة على $7$ بدون باقٍ، ويكرّر الطلب حتى يُدخل المستخدم عددًا يقبل القسمة على $7$ ثم يطبعه:
\begin{boxCode}
\begin{english}
\begin{minted}{csharp}
int n = int.Parse(Console.ReadLine());
while ( ____________ )
{
    n = ____________________________;
}
Console.WriteLine( ________________ );
\end{minted}
\end{english}
\end{boxCode}

    \item البرنامج التالي يستقبل كلمة المرور (عدد صحيح)، وطالما كانت غير مطابقة لكلمة المرور الصحيحة $8264$ يطلب إدخالها من جديد، وفي النهاية يطبع \textenglish{"Welcome"}:
\begin{boxCode}
\begin{english}
\begin{minted}{csharp}
int pass = int.Parse(Console.ReadLine());
while ( ____________ )
{
    Console.WriteLine("Wrong password, try again:");
    ____________________________________;
}
Console.WriteLine("Welcome");
\end{minted}
\end{english}
\end{boxCode}
\end{enumerate}
\ifwithsols
\begin{boxSolution}
\begin{english}
\begin{minted}{csharp}
// a
int n = int.Parse(Console.ReadLine());
while (n % 7 != 0)
{
    n = int.Parse(Console.ReadLine());
}
Console.WriteLine(n);

// b
int pass = int.Parse(Console.ReadLine());
while (pass != 8264)
{
    Console.WriteLine("Wrong password, try again:");
    pass = int.Parse(Console.ReadLine());
}
Console.WriteLine("Welcome");
\end{minted}
\end{english}
\end{boxSolution}
\fi
\filbreak

\item \textbf{المستوى $3$ — أكمل حلقات تحتوي على عدّاد (\textenglish{count}) أو مجموع (\textenglish{sum}):}
\begin{enumerate}
    \item البرنامج التالي يستقبل أعدادًا صحيحة حتى يُدخل المستخدم العدد $0$، وفي النهاية يطبع \textbf{عدد} الأعداد التي أُدخلت (بدون حساب الـ $0$):
\begin{boxCode}
\begin{english}
\begin{minted}{csharp}
int count = ____;
int n = int.Parse(Console.ReadLine());
while ( ____________ )
{
    count = ____________;
    n = int.Parse(Console.ReadLine());
}
Console.WriteLine( ________________ );
\end{minted}
\end{english}
\end{boxCode}

    \item البرنامج التالي يستقبل أسعار كتب (أعدادًا صحيحة) ويجمعها، ويتوقف عندما يُدخل المستخدم السعر $0$، ثم يطبع المجموع النهائي:
\begin{boxCode}
\begin{english}
\begin{minted}{csharp}
int sum = ____;
int price = int.Parse(Console.ReadLine());
while ( ____________ )
{
    sum = ____________;
    price = int.Parse(Console.ReadLine());
}
Console.WriteLine( ________________ );
\end{minted}
\end{english}
\end{boxCode}

    \item البرنامج التالي يستقبل أعدادًا صحيحة حتى يُدخل المستخدم $-1$، ويعدّ ويطبع في النهاية كم عددًا من بينها كان \textbf{موجبًا} (أكبر من $0$):
\begin{boxCode}
\begin{english}
\begin{minted}{csharp}
int posCount = 0;
int n = int.Parse(Console.ReadLine());
while ( ____________ )
{
    if ( ____________ )
    {
        posCount++;
    }
    n = ________________________________;
}
Console.WriteLine( ________________ );
\end{minted}
\end{english}
\end{boxCode}
\end{enumerate}
\ifwithsols
\begin{boxSolution}
\begin{english}
\begin{minted}{csharp}
// a
int count = 0;
int n = int.Parse(Console.ReadLine());
while (n != 0)
{
    count = count + 1; // or count++;
    n = int.Parse(Console.ReadLine());
}
Console.WriteLine(count);

// b
int sum = 0;
int price = int.Parse(Console.ReadLine());
while (price != 0)
{
    sum = sum + price; // or sum += price;
    price = int.Parse(Console.ReadLine());
}
Console.WriteLine(sum);

// c
int posCount = 0;
int n = int.Parse(Console.ReadLine());
while (n != -1)
{
    if (n > 0)
    {
        posCount++;
    }
    n = int.Parse(Console.ReadLine());
}
Console.WriteLine(posCount);
\end{minted}
\end{english}
\end{boxSolution}
\fi
\filbreak

\end{enumerate}

\clearpage
% ==============================================================
% القسم الثاني: أسئلة تأسيسية على حلقة while
% ==============================================================
\section*{القسم الثاني: أسئلة تأسيسية على حلقة \textenglish{while}}

\begin{enumerate}[resume, itemsep=1em]

\item اكتب برنامجًا يستقبل من المستخدم أعدادًا صحيحة، ويطبع تربيع كل عدد يُدخله المستخدم ($n \times n$). إذا أدخل المستخدم العدد $-1$ يتوقف البرنامج (دون طباعة تربيع الـ $-1$).
\ifwithsols
\begin{boxSolution}
\begin{english}
\begin{minted}{csharp}
int n = int.Parse(Console.ReadLine());
while (n != -1)
{
    Console.WriteLine(n * n);
    n = int.Parse(Console.ReadLine());
}
\end{minted}
\end{english}
\end{boxSolution}
\fi
\filbreak

\item اكتب برنامجًا يستقبل من المستخدم أعدادًا صحيحة، ويطبع فقط الأعداد الموجبة من بينها (الأكبر من $0$). ينتهي الاستقبال عندما يُدخل المستخدم العدد $-999$.
\ifwithsols
\begin{boxSolution}
\begin{english}
\begin{minted}{csharp}
int n = int.Parse(Console.ReadLine());
while (n != -999)
{
    if (n > 0)
    {
        Console.WriteLine(n);
    }
    n = int.Parse(Console.ReadLine());
}
\end{minted}
\end{english}
\end{boxSolution}
\fi
\filbreak

\item اكتب برنامجًا يطلب من المستخدم إدخال عدد صحيح زوجي. طالما يُدخل المستخدم عددًا فرديًا، يطبع البرنامج الرسالة \textenglish{"Odd! Try again"} ويطلب إدخال عدد جديد، وعندما يُدخل المستخدم عددًا زوجيًا يطبعه ويتوقف.
\ifwithsols
\begin{boxSolution}
\begin{english}
\begin{minted}{csharp}
int n = int.Parse(Console.ReadLine());
while (n % 2 != 0)
{
    Console.WriteLine("Odd! Try again");
    n = int.Parse(Console.ReadLine());
}
Console.WriteLine(n);
\end{minted}
\end{english}
\end{boxSolution}
\fi
\filbreak

\item اكتب برنامجًا يطلب من المستخدم إدخال عدد صحيح موجب مكوّن من منزلتين (بين $10$ و $99$). طالما أنّ العدد المدخل خارج هذا المجال، يطبع البرنامج \textenglish{"Invalid"} ويطلب إدخال عدد جديد. عندما يُدخل المستخدم عددًا قانونيًا يطبع البرنامج \textenglish{"Valid: "} متبوعًا بالعدد.
\ifwithsols
\begin{boxSolution}
\begin{english}
\begin{minted}{csharp}
int num = int.Parse(Console.ReadLine());
while (num < 10 || num > 99)
{
    Console.WriteLine("Invalid");
    num = int.Parse(Console.ReadLine());
}
Console.WriteLine("Valid: " + num);
\end{minted}
\end{english}
\end{boxSolution}
\fi
\filbreak

\item اكتب برنامجًا يستقبل علامات طلاب في موضوع الرياضيات، ويطبع بالنسبة لكل علامة:
\begin{itemize}[nosep]
    \item \textenglish{"Excellent!"} إذا كانت العلامة أكبر من $90$.
    \item \textenglish{"Nice Work!"} إذا كانت العلامة أكبر من $80$ وحتى $90$.
    \item \textenglish{"Good!"} إذا كانت العلامة أكبر من $60$ وحتى $80$.
    \item \textenglish{"Failed!"} في كل حالة أخرى.
\end{itemize}
ينتهي الاستقبال عندما يُدخل المستخدم علامة غير قانونية (العلامة القانونية هي بين $0$ و $100$).
\ifwithsols
\begin{boxSolution}
\begin{english}
\begin{minted}{csharp}
int grade = int.Parse(Console.ReadLine());
while (grade >= 0 && grade <= 100)
{
    if (grade > 90)
    {
        Console.WriteLine("Excellent!");
    }
    else if (grade > 80)
    {
        Console.WriteLine("Nice Work!");
    }
    else if (grade > 60)
    {
        Console.WriteLine("Good!");
    }
    else
    {
        Console.WriteLine("Failed!");
    }
    grade = int.Parse(Console.ReadLine());
}
\end{minted}
\end{english}
\end{boxSolution}
\fi
\filbreak

\item أمامك مقطع البرنامج التالي:
\begin{boxCode}
\begin{english}
\begin{minted}{csharp}
while (x != y)
{
    x += 1;
    y -= 2;
}
\end{minted}
\end{english}
\end{boxCode}
\begin{enumerate}
    \item أعطِ مثالاً لعددين $x$ و $y$ بحيث تتكرّر الحلقة $3$ مرّات بالضبط.
    \item أعطِ مثالاً لعددين $x$ و $y$ بحيث تتكرّر الحلقة مرّة واحدة فقط.
    \item أعطِ مثالاً لعددين $x$ و $y$ بحيث لا تتكرّر الحلقة أبدًا ($0$ مرّات).
    \item أعطِ مثالاً لعددين $x$ و $y$ بحيث تتكرّر الحلقة إلى ما لا نهاية من المرّات.
\end{enumerate}
\ifwithsols
\begin{boxSolution}
\begin{enumerate}
    \item نأخذ مثلاً: $x = 0$ و $y = 9$. بعد $3$ خطوات يصبحان متساويين ($3$ و $3$) فتتوقف الحلقة (أو أي عددين يحققان $y - x = 9$).
    \item نأخذ مثلاً: $x = 0$ و $y = 3$. بعد خطوة واحدة يصبحان متساويين ($1$ و $1$) فتتوقف الحلقة (أو أي عددين يحققان $y - x = 3$).
    \item أي عددين متساويين منذ البداية (مثلاً $x = 5$ و $y = 5$).
    \item نأخذ مثلاً: $x = 1$ و $y = 0$؛ عندها لن يتساويا أبدًا (أو أي قيمتين لا يقبل فيهما الفرق $y - x$ القسمة على $3$).
\end{enumerate}
\end{boxSolution}
\fi
\filbreak

\end{enumerate}

\clearpage
% ==============================================================
% القسم الثالث: أسئلة تأسيسية على حلقة for
% ==============================================================
\section*{القسم الثالث: أسئلة تأسيسية على حلقة \textenglish{for}}

\begin{enumerate}[resume, itemsep=1em]

\item سجّل المخرج الذي نحصل عليه بالنسبة لكل مقطع من المقاطع التالية (لا حاجة لبناء جدول متابعة):

{\centering
\renewcommand{\arraystretch}{1.3}
\begin{tabular}{|C{1cm}|p{6.2cm}|C{1cm}|p{6.2cm}|}
\hline
$\mathbf{A}$ &
\begin{pycell}
\begin{english}
\begin{minted}{csharp}
for (int i = 1; i <= 5; i++)
{
    Console.WriteLine("!");
}
\end{minted}
\end{english}
\end{pycell}
& $\mathbf{D}$ &
\begin{pycell}
\begin{english}
\begin{minted}{csharp}
for (int i = 1; i <= 10; i += 2)
{
    Console.WriteLine(i * 3);
}
\end{minted}
\end{english}
\end{pycell}
\\ \hline
$\mathbf{B}$ &
\begin{pycell}
\begin{english}
\begin{minted}{csharp}
for (int i = 10; i >= 5; i--)
{
    Console.WriteLine("*");
}
\end{minted}
\end{english}
\end{pycell}
& $\mathbf{E}$ &
\begin{pycell}
\begin{english}
\begin{minted}{csharp}
for (int i = 7; i >= 0; i--)
{
    Console.WriteLine(i * 2);
}
\end{minted}
\end{english}
\end{pycell}
\\ \hline
$\mathbf{C}$ &
\begin{pycell}
\begin{english}
\begin{minted}{csharp}
for (int i = 1; i <= 4; i++)
{
    Console.WriteLine(i);
}
\end{minted}
\end{english}
\end{pycell}
& $\mathbf{F}$ &
\begin{pycell}
\begin{english}
\begin{minted}{csharp}
for (int i = 6; i <= 30; i += 6)
{
    Console.WriteLine(i);
}
\end{minted}
\end{english}
\end{pycell}
\\ \hline
\end{tabular}
\par}
\ifwithsols
\begin{boxSolution}
\begin{itemize}[nosep]
    \item $\mathbf{A}$: يطبع الرمز \textenglish{"!"} خمس مرّات (كل مرّة في سطر).
    \item $\mathbf{B}$: يطبع الرمز \textenglish{"*"} ستّ مرّات (من $10$ إلى $5$).
    \item $\mathbf{C}$: يطبع الأعداد: $1, 2, 3, 4$.
    \item $\mathbf{D}$: يطبع الأعداد: $3, 9, 15, 21, 27$ (لأنّ $i$ يأخذ القيم $1, 3, 5, 7, 9$).
    \item $\mathbf{E}$: يطبع الأعداد: $14, 12, 10, 8, 6, 4, 2, 0$.
    \item $\mathbf{F}$: يطبع الأعداد: $6, 12, 18, 24, 30$.
\end{itemize}
\end{boxSolution}
\fi
\filbreak

\item اكتب برنامجًا يطبع الأعداد من $1$ إلى $20$ باستخدام حلقة \textenglish{for}.
\ifwithsols
\begin{boxSolution}
\begin{english}
\begin{minted}{csharp}
for (int i = 1; i <= 20; i++)
{
    Console.WriteLine(i);
}
\end{minted}
\end{english}
\end{boxSolution}
\fi
\filbreak

\item اكتب برنامجًا يطبع الأعداد من $10$ إلى $1$ تنازليًّا.
\ifwithsols
\begin{boxSolution}
\begin{english}
\begin{minted}{csharp}
for (int i = 10; i >= 1; i--)
{
    Console.WriteLine(i);
}
\end{minted}
\end{english}
\end{boxSolution}
\fi
\filbreak

\item اكتب برنامجًا يستقبل من المستخدم عددًا صحيحًا موجبًا $N$، ثم يطبع $N$ أسطر، في كل سطر الكلمة \textenglish{"Hello"}.
\ifwithsols
\begin{boxSolution}
\begin{english}
\begin{minted}{csharp}
int n = int.Parse(Console.ReadLine());
for (int i = 0; i < n; i++)
{
    Console.WriteLine("Hello");
}
\end{minted}
\end{english}
\end{boxSolution}
\fi
\filbreak

\item اكتب برنامجًا يطبع تربيع الأعداد من $1$ إلى $15$.
\ifwithsols
\begin{boxSolution}
\begin{english}
\begin{minted}{csharp}
for (int i = 1; i <= 15; i++)
{
    Console.WriteLine(i * i);
}
\end{minted}
\end{english}
\end{boxSolution}
\fi
\filbreak

\item اكتب برنامجًا يستقبل من المستخدم عددًا صحيحًا موجبًا $n$، ويطبع الأعداد من $1$ إلى $n$ وبجانب كل عدد تربيعه. \\
مثلاً، إذا أدخل المستخدم العدد $5$، يطبع البرنامج:
\begin{boxCode}
\begin{english}
\begin{minted}{text}
1 1
2 4
3 9
4 16
5 25
\end{minted}
\end{english}
\end{boxCode}
\ifwithsols
\begin{boxSolution}
\begin{english}
\begin{minted}{csharp}
int n = int.Parse(Console.ReadLine());
for (int i = 1; i <= n; i++)
{
    Console.WriteLine(i + " " + (i * i));
}
\end{minted}
\end{english}
\end{boxSolution}
\fi
\filbreak

\end{enumerate}

\clearpage
% ==============================================================
% القسم الرابع: أسئلة حلقات سهلة لتذويت كتابة الحلقات وشروطها ومقدار الخطوة
% ==============================================================
\section*{القسم الرابع: تذويت كتابة الحلقات، شروطها، ومقدار الخطوة}

\begin{enumerate}[resume, itemsep=1em]

\item اكتب برنامجًا يطبع جميع الأعداد الزوجية بين $20$ و $50$ (شاملاً $20$ و $50$).
\ifwithsols
\begin{boxSolution}[1: بزيادة الخطوة بمقدار $2$]
\begin{english}
\begin{minted}{csharp}
for (int i = 20; i <= 50; i += 2)
{
    Console.WriteLine(i);
}
\end{minted}
\end{english}
\end{boxSolution}

\filbreak

\begin{boxSolution}[2: بفحص الشرط داخل الحلقة]
\begin{english}
\begin{minted}{csharp}
for (int i = 20; i <= 50; i++)
{
    if (i % 2 == 0)
    {
        Console.WriteLine(i);
    }
}
\end{minted}
\end{english}
\end{boxSolution}
\fi
\filbreak

\item اكتب برنامجًا يطبع جميع الأعداد التي تقبل القسمة على $5$ بدون باقٍ في المجال من $1$ إلى $100$.
\ifwithsols
\begin{boxSolution}[1]
\begin{english}
\begin{minted}{csharp}
for (int i = 5; i <= 100; i += 5)
{
    Console.WriteLine(i);
}
\end{minted}
\end{english}
\end{boxSolution}

\filbreak

\begin{boxSolution}[2]
\begin{english}
\begin{minted}{csharp}
for (int i = 1; i <= 100; i++)
{
    if (i % 5 == 0)
    {
        Console.WriteLine(i);
    }
}
\end{minted}
\end{english}
\end{boxSolution}
\fi
\filbreak

\item اكتب برنامجًا يطبع الأعداد من $50$ نزولاً إلى $5$ بقفزات مقدارها $5$ (أي: $50, 45, 40, \dots, 5$).
\ifwithsols
\begin{boxSolution}
\begin{english}
\begin{minted}{csharp}
for (int i = 50; i >= 5; i -= 5)
{
    Console.WriteLine(i);
}
\end{minted}
\end{english}
\end{boxSolution}
\fi
\filbreak

\item اكتب برنامجًا يطبع أوّل $12$ مضاعفًا موجبًا للعدد $7$ (أي $7, 14, 21, \dots, 84$).
\ifwithsols
\begin{boxSolution}
\begin{english}
\begin{minted}{csharp}
for (int i = 1; i <= 12; i++)
{
    Console.WriteLine(7 * i);
}
\end{minted}
\end{english}
\end{boxSolution}
\fi
\filbreak

\item اكتب برنامجًا يطبع لكل عدد من $1$ إلى $10$ قيمة العدد مرفوعًا لقوة نفسه \\ (أي: $1^1, 2^2, 3^3, \ldots, 10^{10}$).
\ifwithsols
\begin{boxSolution}
\begin{english}
\begin{minted}{csharp}
for (int i = 1; i <= 10; i++)
{
    Console.WriteLine(Math.Pow(i, i));
}
\end{minted}
\end{english}
\end{boxSolution}
\fi
\filbreak

\item اكتب برنامجًا يستقبل من المستخدم عددًا صحيحًا موجبًا $n$، ويطبع الأعداد من $1$ حتى $n$، ولكن بدلاً من كل عدد يقبل القسمة على $4$ بدون باقٍ يطبع الكلمة \textenglish{"FOUR"} (مثلاً: $1, 2, 3, \text{\textenglish{FOUR}}, 5, \dots$).
\ifwithsols
\begin{boxSolution}
\begin{english}
\begin{minted}{csharp}
int n = int.Parse(Console.ReadLine());
for (int i = 1; i <= n; i++)
{
    if (i % 4 == 0)
    {
        Console.WriteLine("FOUR");
    }
    else
    {
        Console.WriteLine(i);
    }
}
\end{minted}
\end{english}
\end{boxSolution}
\fi
\filbreak

\item اكتب برنامجًا يستقبل $20$ زوجًا من الأعداد الصحيحة، ويطبع بالنسبة لكل زوج العدد الأصغر فيه.
\ifwithsols
\begin{boxSolution}
\begin{english}
\begin{minted}{csharp}
for (int i = 0; i < 20; i++)
{
    int a = int.Parse(Console.ReadLine());
    int b = int.Parse(Console.ReadLine());
    Console.WriteLine(Math.Min(a, b));
}
\end{minted}
\end{english}
\end{boxSolution}
\fi
\filbreak

\item اكتب برنامجًا يستقبل من المستخدم $30$ عددًا صحيحًا موجبًا ثلاثي المنزلة (بين $100$ و $999$)، ويطبع فقط الأعداد التي تتساوى فيها منزلة الآحاد مع منزلة المئات (مثل $141$ و $727$).
\ifwithsols
\begin{boxSolution}
\begin{english}
\begin{minted}{csharp}
for (int i = 0; i < 30; i++)
{
    int num = int.Parse(Console.ReadLine());
    int units = num % 10;
    int hundreds = num / 100;
    if (units == hundreds)
    {
        Console.WriteLine(num);
    }
}
\end{minted}
\end{english}
\end{boxSolution}
\fi
\filbreak

\item اكتب برنامجًا يستقبل أعدادًا صحيحة موجبة ثنائية المنزلة (بين $10$ و $99$)، ويطبع لكل عدد مجموع منزلتيه (الآحاد والعشرات). ينتهي الاستقبال عندما يُدخل المستخدم عددًا ليس ثنائي المنزلة.
\ifwithsols
\begin{boxSolution}
\begin{english}
\begin{minted}{csharp}
int num = int.Parse(Console.ReadLine());
while (num >= 10 && num <= 99)
{
    int sumDigits = (num % 10) + (num / 10);
    Console.WriteLine(sumDigits);
    num = int.Parse(Console.ReadLine());
}
\end{minted}
\end{english}
\end{boxSolution}
\fi
\filbreak

\end{enumerate}

\clearpage
% ==============================================================
% القسم الخامس: أسئلة المجموع (قالب المجموع)
% ==============================================================
\section*{القسم الخامس: أسئلة المجموع (قالب المجموع)}

\begin{boxNote}[- قالب المجموع]
لحساب مجموع قيم في حلقة:
\begin{enumerate}[nosep]
    \item نعرّف متغير المجموع قبل الحلقة ونعطيه القيمة الأولية $0$: \textenglish{int sum = 0;} (وفي حالة الضرب التراكمي نبدأ بـ $1$).
    \item داخل الحلقة نضيف القيمة المطلوبة إلى المجموع: \textenglish{sum += x;}.
    \item بعد انتهاء الحلقة نطبع المجموع النهائي: \textenglish{Console.WriteLine(sum);}.
\end{enumerate}
\end{boxNote}

\begin{enumerate}[resume, itemsep=1em]

\item اكتب برنامجًا يحسب ويطبع مجموع الأعداد من $70$ إلى $100$ (شاملاً $70$ و $100$).
\ifwithsols
\begin{boxSolution}
\begin{english}
\begin{minted}{csharp}
int sum = 0;
for (int i = 70; i <= 100; i++)
{
    sum += i;
}
Console.WriteLine(sum);
\end{minted}
\end{english}
\end{boxSolution}
\fi
\filbreak

\item اكتب برنامجًا يستقبل من المستخدم عددًا صحيحًا موجبًا $n$، ثم يستقبل منه $n$ أعداد صحيحة ويطبع مجموعها.
\ifwithsols
\begin{boxSolution}
\begin{english}
\begin{minted}{csharp}
int n = int.Parse(Console.ReadLine());
int sum = 0;
for (int i = 0; i < n; i++)
{
    int x = int.Parse(Console.ReadLine());
    sum += x;
}
Console.WriteLine(sum);
\end{minted}
\end{english}
\end{boxSolution}
\fi
\filbreak

\item اكتب برنامجًا يستقبل من المستخدم عددين صحيحين $a$ و $b$ (افترض أنّ $a$ أصغر من $b$)، ويحسب ويطبع مجموع جميع الأعداد الصحيحة بينهما (بما يشمل العددين نفسيهما).
\ifwithsols
\begin{boxSolution}
\begin{english}
\begin{minted}{csharp}
int a = int.Parse(Console.ReadLine());
int b = int.Parse(Console.ReadLine());
int sum = 0;
for (int i = a; i <= b; i++)
{
    sum += i;
}
Console.WriteLine(sum);
\end{minted}
\end{english}
\end{boxSolution}
\fi
\filbreak

\item يريد صاحب دكّان صغير حساب مدخوله اليومي؛ اكتب برنامجًا يستقبل المبلغ الذي يدفعه كل زبون عند شرائه من الدكّان، ويحسب ويطبع في النهاية مدخول الدكّان الكلي في ذلك اليوم. ينتهي الاستقبال عند إدخال العدد $-1$.
\ifwithsols
\begin{boxSolution}
\begin{english}
\begin{minted}{csharp}
double sum = 0;
double price = double.Parse(Console.ReadLine());
while (price != -1)
{
    sum += price;
    price = double.Parse(Console.ReadLine());
}
Console.WriteLine(sum);
\end{minted}
\end{english}
\end{boxSolution}
\fi
\filbreak

\item اكتب برنامجًا يستقبل $30$ عددًا صحيحًا، ويحسب ويطبع \textbf{مجموع} الأعداد الزوجية الموجبة المكوّنة من منزلتين فقط (أي الأعداد الزوجية التي بين $10$ و $99$).
\ifwithsols
\begin{boxSolution}
\begin{english}
\begin{minted}{csharp}
int sum = 0;
for (int i = 0; i < 30; i++)
{
    int num = int.Parse(Console.ReadLine());
    if (num >= 10 && num <= 99 && num % 2 == 0)
    {
        sum += num;
    }
}
Console.WriteLine(sum);
\end{minted}
\end{english}
\end{boxSolution}
\fi
\filbreak

\item اكتب برنامجًا يستقبل أعدادًا صحيحة من المستخدم ويجمعها، ويتوقف فقط عندما تصبح قيمة المجموع أكبر تمامًا من $100$، ثم يطبع المجموع النهائي.
\ifwithsols
\begin{boxSolution}
\begin{english}
\begin{minted}{csharp}
int sum = 0;
while (sum <= 100)
{
    int v = int.Parse(Console.ReadLine());
    sum += v;
}
Console.WriteLine(sum);
\end{minted}
\end{english}
\end{boxSolution}
\fi
\filbreak

\item اكتب برنامجًا يستقبل من المستخدم عددًا صحيحًا موجبًا $n$، ويحسب ويطبع مضروب العدد $n$ (الذي يُرمز له بـ $n!$ ويساوي حاصل الضرب $1 \times 2 \times 3 \times \dots \times n$). \\
\textbf{مثال:} إذا أدخل المستخدم العدد $5$، يطبع البرنامج $120$ لأنّ $1 \times 2 \times 3 \times 4 \times 5 = 120$.
\ifwithsols
\begin{boxSolution}
\begin{english}
\begin{minted}{csharp}
int n = int.Parse(Console.ReadLine());
int fact = 1;
for (int i = 1; i <= n; i++)
{
    fact *= i;
}
Console.WriteLine(fact);
\end{minted}
\end{english}
\end{boxSolution}
\fi
\filbreak

\end{enumerate}

\clearpage
% ==============================================================
% القسم السادس: أسئلة العدّ (قالب العدّاد)
% ==============================================================
\section*{القسم السادس: أسئلة العدّ (قالب العدّاد)}

\begin{boxNote}[- قالب العدّاد]
لعدّ المرّات التي يتحقّق فيها شرط معيّن داخل حلقة:
\begin{enumerate}[nosep]
    \item نعرّف متغير العدّاد قبل الحلقة ونعطيه القيمة الأولية $0$: \textenglish{int count = 0;}.
    \item داخل الحلقة نفحص الشرط، وإذا تحقّق نزيد العدّاد بـ $1$: \textenglish{if (...) \{ count++; \}}.
    \item بعد انتهاء الحلقة نطبع العدّاد النهائي: \textenglish{Console.WriteLine(count);}.
\end{enumerate}
\end{boxNote}

\begin{enumerate}[resume, itemsep=1em]

\item وظيفة مقاطع البرامج الأربعة التالية هي استقبال طول $40$ طالبًا (بالأمتار) وطباعة \textbf{عدد} الطلاب الذين يزيد طولهم عن $1.80$ متر. \\
حدّد ما هو المقطع الصحيح، واشرح الخطأ في كل واحد من المقاطع الثلاثة الخاطئة:

{\centering
\renewcommand{\arraystretch}{1.3}
\begin{tabular}{|C{0.8cm}|p{6.6cm}|C{0.8cm}|p{6.6cm}|}
\hline
$\mathbf{A}$ &
\begin{pycell}
\begin{english}
\begin{minted}[fontsize=\footnotesize]{csharp}
for (int i = 0; i < 40; i++)
{
    int count = 0;
    double n = double.Parse(Console.ReadLine());
    if (n > 1.80)
    {
        count++;
    }
    Console.WriteLine(count);
}
\end{minted}
\end{english}
\end{pycell}
& $\mathbf{B}$ &
\begin{pycell}
\begin{english}
\begin{minted}[fontsize=\footnotesize]{csharp}
int count = 0;
for (int i = 1; i <= 40; i++)
{
    double n = double.Parse(Console.ReadLine());
    if (n > 1.80)
    {
        i++;
    }
}
Console.WriteLine(count);
\end{minted}
\end{english}
\end{pycell}
\\ \hline
$\mathbf{C}$ &
\begin{pycell}
\begin{english}
\begin{minted}[fontsize=\footnotesize]{csharp}
int count = 0;
for (int i = 0; i < 40; i++)
{
    double n = double.Parse(Console.ReadLine());
    if (n > 1.80)
    {
        count++;
    }
}
Console.WriteLine(count);
\end{minted}
\end{english}
\end{pycell}
& $\mathbf{D}$ &
\begin{pycell}
\begin{english}
\begin{minted}[fontsize=\footnotesize]{csharp}
int count = 0;
double n = double.Parse(Console.ReadLine());
for (int i = 1; i <= 40; i++)
{
    if (n > 1.80)
    {
        count++;
    }
}
Console.WriteLine(count);
\end{minted}
\end{english}
\end{pycell}
\\ \hline
\end{tabular}
\par}
\ifwithsols
\begin{boxSolution}
\begin{itemize}[nosep]
    \item \textbf{المقطع الصحيح هو} $\mathbf{C}$.
    \item \textbf{الخطأ في} $\mathbf{A}$: تم تصفير \textenglish{count = 0} وطباعته داخل الحلقة في كل دورة بدلاً من خارجها.
    \item \textbf{الخطأ في} $\mathbf{B}$: عند تحقّق الشرط زاد عدّاد الحلقة \textenglish{i++} بدلاً من زيادة عدّاد الطلاب \textenglish{count++}.
    \item \textbf{الخطأ في} $\mathbf{D}$: تم استقبال الطول $n$ مرّة واحدة فقط قبل الحلقة بدلاً من استقباله داخل الحلقة لكل طالب من الـ $40$.
\end{itemize}
\end{boxSolution}
\fi
\filbreak

\item اكتب برنامجًا يحسب ويطبع عدد الأعداد ثلاثية المنزلة (بين $100$ و $999$) التي تقبل القسمة على $7$ بدون باقٍ.
\ifwithsols
\begin{boxSolution}
\begin{english}
\begin{minted}{csharp}
int count = 0;
for (int i = 100; i <= 999; i++)
{
    if (i % 7 == 0)
    {
        count++;
    }
}
Console.WriteLine(count);
\end{minted}
\end{english}
\end{boxSolution}
\fi
\filbreak

\item اكتب برنامجًا يستقبل من المستخدم عددين صحيحين $a$ و $b$ (افترض أنّ العدد الأوّل أصغر من الثاني)، ويحسب ويطبع عدد الأعداد التي تقبل القسمة على $3$ بدون باقٍ بينهما (بما يشمل العددين نفسيهما).
\ifwithsols
\begin{boxSolution}
\begin{english}
\begin{minted}{csharp}
int a = int.Parse(Console.ReadLine());
int b = int.Parse(Console.ReadLine());
int count = 0;
for (int i = a; i <= b; i++)
{
    if (i % 3 == 0)
    {
        count++;
    }
}
Console.WriteLine(count);
\end{minted}
\end{english}
\end{boxSolution}
\fi
\filbreak

\item اكتب برنامجًا يطلب من المستخدم إدخال عدد صحيح بين $1$ و $10$. إذا أدخل عددًا خارج هذا المجال، يطبع البرنامج رسالة خطأ \textenglish{"Error"} ويطلب عددًا جديدًا حتى يُدخل عددًا ضمن المجال المطلوب. في النهاية يطبع البرنامج عدد المرات التي أدخل فيها المستخدم عددًا خاطئًا.
\ifwithsols
\begin{boxSolution}
\begin{english}
\begin{minted}{csharp}
int invalid = 0;
int x = int.Parse(Console.ReadLine());
while (x < 1 || x > 10)
{
    Console.WriteLine("Error");
    invalid++;
    x = int.Parse(Console.ReadLine());
}
Console.WriteLine(invalid);
\end{minted}
\end{english}
\end{boxSolution}
\fi
\filbreak

\item اكتب برنامجًا يستقبل في البداية عددًا صحيحًا موجبًا $x$، ثم يستقبل أعدادًا صحيحة موجبة حتى يُدخل المستخدم العدد $0$. في النهاية يطبع البرنامج \textbf{عدد} الأعداد المدخلة التي تقبل القسمة على $x$ بدون باقٍ.
\ifwithsols
\begin{boxSolution}
\begin{english}
\begin{minted}{csharp}
int x = int.Parse(Console.ReadLine());
int count = 0;
int num = int.Parse(Console.ReadLine());
while (num != 0)
{
    if (num % x == 0)
    {
        count++;
    }
    num = int.Parse(Console.ReadLine());
}
Console.WriteLine(count);
\end{minted}
\end{english}
\end{boxSolution}
\fi
\filbreak

\item اكتب برنامجًا يستقبل علامات طلاب (أعدادًا صحيحة بين $0$ و $100$)، ويتوقف عن الاستقبال عندما يُدخل المستخدم العدد $-1$. بعد التوقف، يطبع البرنامج \textbf{عدد} العلامات التي أُدخلت و\textbf{معدّلها} (افترض أنّ المستخدم يُدخل علامة واحدة على الأقل قبل $-1$).
\ifwithsols
\begin{boxSolution}
\begin{english}
\begin{minted}{csharp}
int sum = 0;
int count = 0;
int grade = int.Parse(Console.ReadLine());
while (grade != -1)
{
    sum += grade;
    count++;
    grade = int.Parse(Console.ReadLine());
}
double avg = (double)sum / count;
Console.WriteLine("Count: " + count);
Console.WriteLine("Average: " + avg);
\end{minted}
\end{english}
\end{boxSolution}
\fi
\filbreak

\end{enumerate}

\clearpage
% ==============================================================
% القسم السابع: أسئلة القيمة القصوى والدنيا (max و min)
% ==============================================================
\section*{القسم السابع: أسئلة القيمة القصوى والدنيا (\textenglish{max} و \textenglish{min})}

\begin{boxNote}[- قالب إيجاد القيمة القصوى والدنيا]
\begin{itemize}[nosep]
    \item لإيجاد \textbf{أكبر قيمة (\textenglish{max}):} نعرّف قبل الحلقة \textenglish{int max = int.MinValue;} (أو نعطيه أوّل قيمة مدخلة)، وداخل الحلقة نفحص: \textenglish{if (x > max) \{ max = x; \}}.
    \item لإيجاد \textbf{أصغر قيمة (\textenglish{min}):} نعرّف قبل الحلقة \textenglish{int min = int.MaxValue;} (أو نعطيه أوّل قيمة مدخلة)، وداخل الحلقة نفحص: \textenglish{if (x < min) \{ min = x; \}}.
\end{itemize}
\end{boxNote}

\begin{enumerate}[resume, itemsep=1em]

\item في كيس هناك $33$ طابة. اكتب برنامجًا يستقبل وزن كل طابة، وفي النهاية يطبع وزن الطابة \textbf{الأكبر}.
\ifwithsols
\begin{boxSolution}
\begin{english}
\begin{minted}{csharp}
double maxWeight = double.MinValue;
for (int i = 0; i < 33; i++)
{
    double w = double.Parse(Console.ReadLine());
    if (w > maxWeight)
    {
        maxWeight = w;
    }
}
Console.WriteLine(maxWeight);
\end{minted}
\end{english}
\end{boxSolution}
\fi
\filbreak

\item أعلنت البلدية عن مناقصة لبناء مدرسة إعدادية جديدة، وتقدّم للمناقصة $15$ مقاول بناء. اكتب برنامجًا يستقبل المبلغ الذي اقترحه كل مقاول للمناقصة، ويطبع في النهاية المبلغ الفائز بالمناقصة وهو \textbf{أقل اقتراح}.
\ifwithsols
\begin{boxSolution}
\begin{english}
\begin{minted}{csharp}
int minOffer = int.MaxValue;
for (int i = 0; i < 15; i++)
{
    int offer = int.Parse(Console.ReadLine());
    if (offer < minOffer)
    {
        minOffer = offer;
    }
}
Console.WriteLine(minOffer);
\end{minted}
\end{english}
\end{boxSolution}
\fi
\filbreak

\item اكتب برنامجًا يستقبل من المستخدم أعدادًا صحيحة حتى يُدخل العدد $999$. بعد التوقف، يطبع البرنامج \textbf{أكبر} عدد قام المستخدم بإدخاله قبل العدد $999$.
\ifwithsols
\begin{boxSolution}
\begin{english}
\begin{minted}{csharp}
int max = int.MinValue;
int v = int.Parse(Console.ReadLine());
while (v != 999)
{
    if (v > max)
    {
        max = v;
    }
    v = int.Parse(Console.ReadLine());
}
Console.WriteLine(max);
\end{minted}
\end{english}
\end{boxSolution}
\fi
\filbreak

\item اكتب برنامجًا يستقبل من المستخدم $100$ عدد صحيح، ويحسب ويطبع \textbf{أكبر عدد زوجي} و\textbf{أكبر عدد فردي} من بين الأعداد التي أدخلها المستخدم.
\ifwithsols
\begin{boxSolution}
\begin{english}
\begin{minted}{csharp}
int maxEven = int.MinValue;
int maxOdd = int.MinValue;
for (int i = 1; i <= 100; i++)
{
    int x = int.Parse(Console.ReadLine());
    if (x % 2 == 0)
    {
        maxEven = Math.Max(maxEven, x);
    }
    else
    {
        maxOdd = Math.Max(maxOdd, x);
    }
}
Console.WriteLine(maxEven);
Console.WriteLine(maxOdd);
\end{minted}
\end{english}
\end{boxSolution}
\fi
\filbreak

\item في إحدى مسابقات ألعاب القوى هناك مسابقة «العرض الشيّق»، وفيها يعطي $7$ حكّام علامة لكل متسابق بعد عرضه. \\
العلامة النهائية التي يستحقها المتسابق تُحسب كمعدّل لـ $5$ علامات فقط من بين العلامات الـ $7$، وذلك بعد شطب العلامة الكبرى والعلامة الصغرى. \\
اكتب برنامجًا يستقبل علامات الحكّام السبعة لأحد المتسابقين، ويطبع معدّله النهائي.
\ifwithsols
\begin{boxSolution}
\begin{english}
\begin{minted}{csharp}
int min = int.MaxValue;
int max = int.MinValue;
int sum = 0;
for (int i = 0; i < 7; i++)
{
    int s = int.Parse(Console.ReadLine());
    sum += s;
    if (s < min)
    {
        min = s;
    }
    if (s > max)
    {
        max = s;
    }
}
double avg = (sum - min - max) / 5.0;
Console.WriteLine(avg);
\end{minted}
\end{english}
\end{boxSolution}
\fi
\filbreak

\end{enumerate}

\clearpage
% ==============================================================
% القسم الثامن: أسئلة حلقات متنوّعة بمستوى متقدّم
% ==============================================================
\section*{القسم الثامن: أسئلة حلقات متنوّعة بمستوى متقدّم}

\begin{enumerate}[resume, itemsep=2em]

\item اكتب برنامجًا يستقبل من المستخدم عددين صحيحين، ويحسب ويطبع مجموع جميع الأعداد الصحيحة بينهما (بما يشمل العددين نفسيهما). \\
\textbf{ملاحظة:} لا تفترض أنّ العدد الأوّل أصغر من الثاني (قد يكون الأوّل أكبر وقد يكون الثاني أكبر).
\ifwithsols
\begin{boxSolution}
\begin{english}
\begin{minted}{csharp}
int a = int.Parse(Console.ReadLine());
int b = int.Parse(Console.ReadLine());
int max = Math.Max(a, b);
int min = Math.Min(a, b);
int sum = 0;
for (int i = min; i <= max; i++)
{
    sum += i;
}
Console.WriteLine(sum);
\end{minted}
\end{english}
\end{boxSolution}
\fi
\filbreak

\item اكتب برنامجًا يستقبل عددًا صحيحًا موجبًا $n$، ويحسب ويطبع قيمة التعبير التالي: \\
$1 - 2 + 3 - 4 + 5 - 6 + \dots \pm n$ \\
(لاحظ أنّ الأعداد الفردية تُضاف، والأعداد الزوجية تُطرح.)
\ifwithsols
\begin{boxSolution}
\begin{english}
\begin{minted}{csharp}
int n = int.Parse(Console.ReadLine());
int sum = 0;
for (int i = 1; i <= n; i++)
{
    if (i % 2 == 0)
    {
        sum -= i;
    }
    else
    {
        sum += i;
    }
}
Console.WriteLine(sum);
\end{minted}
\end{english}
\end{boxSolution}
\fi
\filbreak

\item اكتب برنامجًا يستقبل عددًا صحيحًا موجبًا $n$، ويحسب ويطبع كم عددًا $i$ في المجال من $1$ حتى $n$ يحقّق أنّ قيمة التعبير $i^2 - i$ تقبل القسمة على $3$ بدون باقٍ. \\
(مثلاً: عندما $i = 2$ فإنّ $2^2 - 2 = 2$ لا يقبل القسمة على $3$، أمّا عندما $i = 3$ فإنّ $3^2 - 3 = 6$ يقبل القسمة على $3$ فنعدّه.)
\ifwithsols
\begin{boxSolution}
\begin{english}
\begin{minted}{csharp}
int n = int.Parse(Console.ReadLine());
int cnt = 0;
for (int i = 1; i <= n; i++)
{
    int v = i * i - i;
    if (v % 3 == 0)
    {
        cnt++;
    }
}
Console.WriteLine(cnt);
\end{minted}
\end{english}
\end{boxSolution}
\fi
\filbreak

\item اكتب برنامجًا يستقبل من المستخدم $100$ عدد صحيح، ويعدّ الأعداد التي لها نفس زوجية ترتيب إدخالها ويطبع عددها. \\
أي: إذا كان العدد الأوّل فرديًّا نعدّه وإلا فلا، وإذا كان العدد الثاني زوجيًّا نعدّه وإلا فلا، وإذا كان العدد الثالث فرديًّا نعدّه وإلا فلا، وهكذا.
\ifwithsols
\begin{boxSolution}
\begin{english}
\begin{minted}{csharp}
int cnt = 0;
for (int i = 1; i <= 100; i++)
{
    int x = int.Parse(Console.ReadLine());
    if (i % 2 == Math.Abs(x % 2))
    {
        cnt++;
    }
}
Console.WriteLine(cnt);
\end{minted}
\end{english}
\end{boxSolution}
\fi
\filbreak

\item تُحسب العلامة النهائية في موضوع الحاسوب على النحو التالي:
\begin{itemize}[nosep]
    \item $30\%$ للوظائف
    \item $30\%$ للامتحان الشهري
    \item $40\%$ للامتحان الفصلي
\end{itemize}
اكتب برنامجًا يستقبل عدد طلاب الصف، ثم يستقبل علامات كل واحد من الطلاب (لكل طالب يستقبل $3$ علامات)، ويطبع لكل طالب علامته النهائية. وفي النهاية يطبع البرنامج عدد الطلاب الناجحين الذين حصلوا على علامة نهائية $70$ فما فوق.
\ifwithsols
\begin{boxSolution}
\begin{english}
\begin{minted}{csharp}
int pupilsCnt = int.Parse(Console.ReadLine());
int cnt = 0;
for (int i = 1; i <= pupilsCnt; i++)
{
    int hw = int.Parse(Console.ReadLine());
    int mid = int.Parse(Console.ReadLine());
    int finalExam = int.Parse(Console.ReadLine());
    double avg = hw * 0.3 + mid * 0.3 + finalExam * 0.4;
    Console.WriteLine(avg);
    if (avg >= 70)
    {
        cnt++;
    }
}
Console.WriteLine(cnt);
\end{minted}
\end{english}
\end{boxSolution}
\fi
\filbreak

\item اكتب برنامجًا يستقبل $25$ عددًا صحيحًا، ويعدّ ويطبع كم عددًا منها كان أصغر من معدّل الأعداد التي أُدخلت منذ البداية وحتى تلك اللحظة (أي في الخطوة $i$، نحسب معدّل أوّل $i$ أعداد، ونفحص هل العدد الـ $i$ أصغر من هذا المعدّل).
\ifdetailed
\begin{boxExample}
إذا كانت أوّل $5$ أعداد مدخلة هي: $12, 8, 20, 4, 6$ فإنّ البرنامج يحسب كالتالي:
\begin{enumerate}[label=(\arabic*), nosep]
    \item عند قراءة العدد $12$: المتوسط الجزئي هو $12$. العدد $12$ ليس أصغر من $12$ $\leftarrow$ لا نعدّه.
    \item عند قراءة العدد $8$: المتوسط الجزئي هو $\frac{12 + 8}{2} = 10$. العدد $8$ أصغر من $10$ $\leftarrow$ نعدّه (العداد الآن $1$).
    \item عند قراءة العدد $20$: المتوسط الجزئي هو $\frac{12 + 8 + 20}{3} = 13.33$. العدد $20$ ليس أصغر من $13.33$ $\leftarrow$ لا نعدّه (العداد الآن $1$).
    \item عند قراءة العدد $4$: المتوسط الجزئي هو $\frac{12 + 8 + 20 + 4}{4} = 11$. العدد $4$ أصغر من $11$ $\leftarrow$ نعدّه (العداد الآن $2$).
    \item عند قراءة العدد $6$: المتوسط الجزئي هو $\frac{12 + 8 + 20 + 4 + 6}{5} = 10$. العدد $6$ أصغر من $10$ $\leftarrow$ نعدّه (العداد الآن $3$).
\end{enumerate}
\end{boxExample}
\fi
\ifwithsols
\begin{boxSolution}
\begin{english}
\begin{minted}{csharp}
int sum = 0;
int cnt = 0;
for (int i = 1; i <= 25; i++)
{
    int x = int.Parse(Console.ReadLine());
    sum += x;
    double avg = (double)sum / i;
    if (x < avg)
    {
        cnt++;
    }
}
Console.WriteLine(cnt);
\end{minted}
\end{english}
\end{boxSolution}
\fi
\filbreak

\item اكتب برنامجًا يستقبل من المستخدم $20$ زوجًا من الأعداد الصحيحة (أي ما مجمله $40$ عددًا). \\
لكل زوج من الأعداد، يحسب البرنامج معدّل العددين ويطبعه، وفي النهاية يطبع العددين اللذين كان لهما أكبر معدّل.
\ifdetailed
\begin{boxExample}
لو كان البرنامج يستقبل $5$ أزواج فقط وهي: $(2,4), (5,7), (1,9), (6,8), (3,3)$ فإنّه يعمل كما يلي:
\begin{enumerate}[label=(\arabic*), nosep]
    \item يستقبل العددين $2$ و $4$، المعدّل هو $\frac{2+4}{2} = 3$، فيطبع $3$.
    \item يستقبل العددين $5$ و $7$، المعدّل هو $\frac{5+7}{2} = 6$، فيطبع $6$.
    \item يستقبل العددين $1$ و $9$، المعدّل هو $\frac{1+9}{2} = 5$، فيطبع $5$.
    \item يستقبل العددين $6$ و $8$، المعدّل هو $\frac{6+8}{2} = 7$، فيطبع $7$.
    \item يستقبل العددين $3$ و $3$، المعدّل هو $\frac{3+3}{2} = 3$، فيطبع $3$.
\end{enumerate}
وفي النهاية يطبع العددين اللذين كان لهما أكبر معدّل وهما: $6$ و $8$ (لأنّ معدّلهما $7$ هو الأكبر).
\end{boxExample}
\fi
\ifwithsols
\begin{boxSolution}
\begin{english}
\begin{minted}{csharp}
double maxAvg = double.MinValue;
int num1 = 0, num2 = 0;
for (int i = 0; i < 20; i++)
{
    int a = int.Parse(Console.ReadLine());
    int b = int.Parse(Console.ReadLine());
    double avg = (a + b) / 2.0;
    Console.WriteLine(avg);
    if (avg > maxAvg)
    {
        maxAvg = avg;
        num1 = a;
        num2 = b;
    }
}
Console.WriteLine(num1);
Console.WriteLine(num2);
\end{minted}
\end{english}
\end{boxSolution}
\fi
\filbreak

\item اكتب برنامجًا يستقبل من المستخدم $20$ عددًا صحيحًا، ويطبع أكبر عددين أدخلهما المستخدم (أكبر عدد وثاني أكبر عدد).
\ifwithsols
\begin{boxSolution}
\begin{english}
\begin{minted}{csharp}
int max1 = int.MinValue;
int max2 = int.MinValue;
for (int i = 0; i < 20; i++)
{
    int x = int.Parse(Console.ReadLine());
    if (x > max1)
    {
        max2 = max1;
        max1 = x;
    }
    else if (x > max2)
    {
        max2 = x;
    }
}
Console.WriteLine(max1);
Console.WriteLine(max2);
\end{minted}
\end{english}
\end{boxSolution}
\fi
\filbreak

\item اكتب برنامجًا يستقبل من المستخدم $20$ عددًا صحيحًا، ويطبع أصغر عددين أدخلهما المستخدم (أصغر عدد وثاني أصغر عدد).
\ifwithsols
\begin{boxSolution}
\begin{english}
\begin{minted}{csharp}
int min1 = int.MaxValue;
int min2 = int.MaxValue;
for (int i = 0; i < 20; i++)
{
    int x = int.Parse(Console.ReadLine());
    if (x < min1)
    {
        min2 = min1;
        min1 = x;
    }
    else if (x < min2)
    {
        min2 = x;
    }
}
Console.WriteLine(min1);
Console.WriteLine(min2);
\end{minted}
\end{english}
\end{boxSolution}
\fi
\filbreak

\item اكتب برنامجًا يستقبل $20$ عددًا صحيحًا من المستخدم، ويطبع عدد المرات التي أدخل فيها المستخدم عددين متتاليين متساويين. \\
\textbf{مثال:} إذا كانت المدخلات $5, 3, \mathbf{4, 4}, 7, \mathbf{9, 9} \dots$ $\leftarrow$ العداد هنا يساوي $2$.
\ifwithsols
\begin{boxSolution}
\begin{english}
\begin{minted}{csharp}
int count = 0;
int prev = int.Parse(Console.ReadLine());
for (int i = 1; i < 20; i++)
{
    int curr = int.Parse(Console.ReadLine());
    if (curr == prev)
    {
        count++;
    }
    prev = curr;
}
Console.WriteLine(count);
\end{minted}
\end{english}
\end{boxSolution}
\fi
\filbreak

\item اكتب برنامجًا يستقبل من المستخدم $10$ أعداد صحيحة، ويطبع عدد الأعداد التي كانت أكبر من العدد الذي أُدخل قبلها مباشرة.
\ifdetailed
\begin{boxExample}
إذا كانت الأعداد العشرة المدخلة هي: $5, 8, 3, 10, 7, 12, 6, 15, 9, 20$:
\begin{itemize}[nosep]
    \item الأعداد التي كانت أكبر من العدد الذي سبقها مباشرة هي: $8$ (أكبر من $5$)، $10$ (أكبر من $3$)، $12$ (أكبر من $7$)، $15$ (أكبر من $6$)، و $20$ (أكبر من $9$).
    \item يطبع البرنامج في النهاية $5$.
\end{itemize}
\end{boxExample}
\fi
\ifwithsols
\begin{boxSolution}
\begin{english}
\begin{minted}{csharp}
int prev = int.Parse(Console.ReadLine());
int count = 0;
for (int i = 1; i < 10; i++)
{
    int current = int.Parse(Console.ReadLine());
    if (current > prev)
    {
        count++;
    }
    prev = current;
}
Console.WriteLine(count);
\end{minted}
\end{english}
\end{boxSolution}
\fi
\filbreak

\item اكتب برنامجًا يستقبل $30$ عددًا صحيحًا، ويحسب ويطبع أكبر فرق (بالقيمة المطلقة) بين عددين متتاليين. \\
(مثلاً: لو كانت الأعداد الخمسة الأولى هي $10, 18, 15, 3, 9$، فإنّ الفروق المتتالية هي $|18 - 10| = 8$، ثم $|15 - 18| = 3$، ثم $|3 - 15| = 12$، ثم $|9 - 3| = 6$، وأكبر فرق بينها هو $12$.)
\ifwithsols
\begin{boxSolution}
\begin{english}
\begin{minted}{csharp}
int prev = int.Parse(Console.ReadLine());
int maxDiff = 0;
for (int i = 1; i < 30; i++)
{
    int x = int.Parse(Console.ReadLine());
    int diff = Math.Abs(x - prev);
    if (diff > maxDiff)
    {
        maxDiff = diff;
    }
    prev = x;
}
Console.WriteLine(maxDiff);
\end{minted}
\end{english}
\end{boxSolution}
\fi
\filbreak

\item اكتب برنامجًا يستقبل من المستخدم أعدادًا صحيحة، ويطبع الجذر التربيعي لكل عدد غير سالب يدخله المستخدم (بعد التأكّد من أنّ العدد ليس سالبًا). ينتهي الاستقبال عندما يصبح عدد الأعداد السالبة التي أدخلها المستخدم $3$ أعداد.
\ifwithsols
\begin{boxSolution}
\begin{english}
\begin{minted}{csharp}
int negCnt = 0;
while (negCnt < 3)
{
    int x = int.Parse(Console.ReadLine());
    if (x >= 0)
    {
        Console.WriteLine(Math.Sqrt(x));
    }
    else
    {
        negCnt++;
    }
}
\end{minted}
\end{english}
\end{boxSolution}
\fi
\filbreak

\item اكتب برنامجًا يستقبل عددًا صحيحًا $x$ يمثّل الحد الأقصى للوزن الإجمالي الذي يمكن أن تحمله شاحنة بالكيلوغرام، ثم يستقبل وزن كل صندوق يتم تحميله للشاحنة الواحد تلو الآخر طالما أنّ الوزن الإجمالي لا يتجاوز $x$. في النهاية يطبع البرنامج عدد الصناديق التي تم تحميلها للشاحنة حتى امتلأت. (فكّر: ما هو شرط الحلقة؟)
\ifwithsols
\begin{boxSolution}
\begin{english}
\begin{minted}{csharp}
int x = int.Parse(Console.ReadLine());
int sum = 0;
int count = 0;
int w = int.Parse(Console.ReadLine());
while (sum + w <= x)
{
    sum += w;
    count++;
    w = int.Parse(Console.ReadLine());
}
Console.WriteLine(count);
\end{minted}
\end{english}
\end{boxSolution}
\fi
\filbreak

\item تنظّم المدرسة حملة تبرّعات، وتستمر الحملة باستقبال مبالغ التبرّعات (أعداد صحيحة موجبة) حتى يتحقّق \textbf{أحد} الشرطين التاليين:
\begin{itemize}[nosep]
    \item أن يصل المجموع الكلي للتبرّعات إلى $5000$ شاقل أو أكثر.
    \item أو أن يصل عدد المتبرّعين إلى $30$ متبرّعًا.
\end{itemize}
اكتب برنامجًا يستقبل مبالغ التبرّعات الواحد تلو الآخر حتى تتوقف الحملة، وفي النهاية يطبع المجموع الكلي للتبرّعات وعدد المتبرّعين.
\ifwithsols
\begin{boxSolution}
\begin{english}
\begin{minted}{csharp}
int totalDonations = 0;
int donorsCount = 0;
while (totalDonations < 5000 && donorsCount < 30)
{
    int amount = int.Parse(Console.ReadLine());
    totalDonations += amount;
    donorsCount++;
}
Console.WriteLine("Total donations: " + totalDonations);
Console.WriteLine("Donors count: " + donorsCount);
\end{minted}
\end{english}
\end{boxSolution}
\fi
\filbreak

\end{enumerate}

\vspace{1cm}
\begin{flushleft}
أرجو لكم وقتًا ممتعًا.

الأستاذ محمود اغبارية.
\end{flushleft}

\end{document}
